\section{Estruturas de dados: da CHE à IBHM}
\label{sec:estruturas}

O algoritmo opera inteiramente sobre a topologia da malha. Esta seção apresenta o vocabulário necessário: o conceito
de \emph{half-edge}, a estrutura compacta CHE, sua extensão para malhas híbridas (IBHM) e como as duas estão
implementadas no repositório.

\subsection{Half-edges}

Em uma malha 2D orientada, cada elemento é um polígono cujos vértices são listados no sentido anti-horário. Cada lado
do elemento, visto com a orientação induzida pelo elemento, é uma \emph{half-edge}~\cite{mantyla}. Uma aresta interna
é, portanto, representada por duas half-edges com sentidos opostos, uma em cada elemento vizinho; uma aresta de borda
tem apenas uma. Cada half-edge $h$ guarda o vértice de onde parte, $\V[h]$. A partir dela definem-se três consultas
(\cref{fig:halfedge}):
\begin{itemize}
\item $\nxt(h)$ e $\prv(h)$: a half-edge seguinte e a anterior no mesmo elemento;
\item $\Op[h]$: a half-edge \emph{oposta}, isto é, a mesma aresta vista pelo elemento vizinho (ou \BORDER{} se $h$
está na borda).
\end{itemize}
Note que $h$ parte de $\V[h]$ e chega em $\V[\nxt(h)]$, e que $\Op[h]$ vai de $\V[\nxt(h)]$ para $\V[h]$.

\begin{figure}[ht]
\centering
\begin{subfigure}[b]{0.44\linewidth}\centering
\begin{tikzpicture}[scale=1.25, every node/.style={font=\small}]
  \filldraw[cefill] (0,0) -- (2,0) -- (2,2) -- (0,2) -- cycle;
  \filldraw[cefill, fill=ceblue!60] (2,0) -- (4,0) -- (4,2) -- (2,2) -- cycle;
  \node[elemlabel] at (0.8,1.05) {$c_0$}; \node[elemlabel] at (3.3,1.05) {$c_1$};
  \draw[hearrow, red!70!black, line width=0.9pt] (1.82,0.35) -- (1.82,1.65) node[midway, left=1pt] {$h$};
  \draw[hearrow, violet] (1.65,1.82) -- (0.35,1.82) node[midway, below=1pt] {$\nxt(h)$};
  \draw[hearrow, green!45!black] (0.35,0.18) -- (1.65,0.18) node[midway, above=1pt] {$\prv(h)$};
  \draw[hearrow, blue!70!black, line width=0.9pt] (2.18,1.65) -- (2.18,0.35) node[midway, right=1pt] {$\Op[h]$};
  \node[vreg, minimum size=0.36cm, label={[font=\scriptsize]below right:{$v=\V[h]$}}] at (2,0) {};
  \node[vreg, minimum size=0.36cm] at (2,2) {}; \node[vreg, dot] at (0,0) {}; \node[vreg, dot] at (0,2) {};
  \node[vreg, dot] at (4,0) {}; \node[vreg, dot] at (4,2) {};
\end{tikzpicture}
\caption{Consultas básicas.}
\label{fig:halfedge}
\end{subfigure}\hfill
\begin{subfigure}[b]{0.52\linewidth}\centering
\begin{tikzpicture}[scale=1.2, every node/.style={font=\small}]
  \foreach \i in {0,...,5} { \coordinate (p\i) at ({60*\i}:2.0); }
  \coordinate (v) at (0,0);
  \foreach \i/\j in {0/1,1/2,2/3,3/4,4/5,5/0} { \filldraw[cefill] (v) -- (p\i) -- (p\j) -- cycle; }
  \filldraw[cefill, fill=ceblue!55] (v) -- (p1) -- (p2) -- cycle;
  \node[elemlabel, font=\scriptsize\sffamily\bfseries] at (90:0.62) {$T_1$};
  \node[elemlabel, font=\scriptsize\sffamily\bfseries] at (150:0.62) {$T_2$};
  \node[elemlabel, font=\scriptsize\sffamily\bfseries] at (30:0.62) {$T_0$};
  % h e prev(h) em T1
  \draw[hearrow, red!70!black, line width=0.9pt] ($(v)!0.2!(p1)+(150:0.1)$) -- ($(v)!0.8!(p1)+(150:0.1)$);
  \node[text=red!70!black] at (73:1.05) {$h$};
  \draw[hearrow, violet] ($(p2)!0.2!(v)+(30:0.1)$) -- ($(p2)!0.8!(v)+(30:0.1)$);
  \node[text=violet, font=\scriptsize] at (98:1.55) {$\prv(h)$};
  % O[prev(h)] em T2
  \draw[hearrow, green!45!black, line width=0.9pt] ($(v)!0.2!(p2)+(210:0.1)$) -- ($(v)!0.8!(p2)+(210:0.1)$);
  \node[text=green!45!black, font=\scriptsize] at (142:1.3) {$\Op[\prv(h)]$};
  % O[h] e next(O[h]) em T0
  \draw[hearrow, blue!70!black] ($(p1)!0.2!(v)+(-30:0.1)$) -- ($(p1)!0.8!(v)+(-30:0.1)$);
  \node[text=blue!70!black, font=\scriptsize] at (44:1.5) {$\Op[h]$};
  \draw[hearrow, orange!80!black, line width=0.9pt] ($(v)!0.2!(p0)+(90:0.1)$) -- ($(v)!0.8!(p0)+(90:0.1)$);
  \node[text=orange!80!black, font=\scriptsize] at (13:1.25) {$\nxt(\Op[h])$};
  \node[vreg, minimum size=0.36cm] at (v) {$v$};
  \foreach \i in {0,...,5} { \node[vreg, dot] at (p\i) {}; }
  \draw[-{Stealth[length=5pt]}, green!45!black, line width=0.8pt] (68:2.3) arc (68:118:2.3)
    node[pos=0.55, above=1pt, font=\scriptsize] {anti-horário};
  \draw[-{Stealth[length=5pt]}, orange!80!black, line width=0.8pt] (52:2.3) arc (52:8:2.3)
    node[midway, right=1pt, font=\scriptsize] {horário};
\end{tikzpicture}
\caption{Girando em torno de um vértice $v$.}
\label{fig:rotacao}
\end{subfigure}
\caption{(a) Uma half-edge $h$ do elemento $c_0$ e suas consultas. (b) As duas maneiras de girar em torno de um vértice
usando apenas $\nxt$, $\prv$ e $\Op$: a regra $h\leftarrow\Op[\prv(h)]$ passa ao elemento vizinho no sentido
anti-horário (usada por \code{reindexElementsCCW}), e a regra $h\leftarrow\nxt(\Op[h])$ passa ao vizinho no sentido
horário (usada por \code{isBorder} e pelo segundo laço de \code{openHole}). Em ambos os casos a nova half-edge também
parte de $v$.}
\end{figure}

\subsection{CHE: half-edges implícitas}

A \emph{Compact Half-Edge} (CHE)~\cite{che} dispensa objetos ``half-edge''. Se todos os elementos têm $k$ vértices,
basta guardar a lista de elementos como um vetor de inteiros $\V[\,]$ com $k$ entradas por elemento: a posição $h$ desse
vetor \emph{é} a half-edge, e o valor armazenado é o vértice de origem. As demais consultas saem por aritmética:
\[
  \mathrm{elem}(h)=\lfloor h/k\rfloor,\qquad
  \nxt(h)=k\,\mathrm{elem}(h)+(h+1)\bmod k,\qquad
  \prv(h)=k\,\mathrm{elem}(h)+(h+k-1)\bmod k .
\]
O único vetor adicional é o das opostas, $\Op[\,]$, com uma entrada por half-edge. Toda a adjacência da malha cabe em
dois vetores de inteiros do mesmo tamanho.

\subsection{IBHM: malhas híbridas}

A IBHM (\emph{Index-Based data structure for Hybrid Meshes})~\cite{ibhm} estende a CHE (e a CHF~\cite{chf}, em 3D) para
malhas com elementos de tipos diferentes. Em 2D, a mudança essencial é um vetor de \emph{offsets} $Of[\,]$ que diz onde
cada elemento começa em $\V[\,]$. Assim, o elemento $c$ ocupa as posições $Of[c],\dots,Of[c]+n_v(c)-1$, e:
\begin{itemize}
\item $\mathrm{elem}(h)$ é obtido por \emph{busca binária} em $Of[\,]$, em $O(\log n_c)$ (em vez de uma divisão);
\item $\nxt(h)=Of[c]+\big((h-Of[c]+1)\bmod n_v(c)\big)$ e analogamente para $\prv$;
\item $\Op[\,]$ é idêntico ao da CHE;
\item um vetor extra $VH[\,]$ associa a cada vértice uma half-edge que parte dele (a de borda, se o vértice estiver na
borda), o que permite iniciar consultas de vizinhança e testar borda em $O(1)$.
\end{itemize}
A \cref{fig:ibhm} mostra os vetores para a malha do arquivo \code{meshIBHM.txt}, que reproduz (com índices a partir de
zero) o exemplo da Figura~4 do artigo: um quadrado com um furo quadrado, formado por três quadriláteros e três
triângulos.

\begin{figure}[ht]
\centering
\begin{tikzpicture}[scale=2.500, every node/.style={font=\scriptsize}]
\fill[holefill] (2.000,2.000) -- (2.000,1.000) -- (1.000,1.000) -- (1.000,2.000) -- cycle;
\filldraw[cefill] (2.000,2.000) -- (3.000,3.000) -- (1.500,3.000) -- cycle;
\filldraw[cefill] (3.000,3.000) -- (2.000,2.000) -- (2.000,1.000) -- (3.000,0.000) -- cycle;
\filldraw[cefill] (3.000,0.000) -- (2.000,1.000) -- (1.000,1.000) -- (0.000,0.000) -- cycle;
\filldraw[cefill] (0.000,0.000) -- (1.000,1.000) -- (1.000,2.000) -- (0.000,3.000) -- cycle;
\filldraw[cefill] (0.000,3.000) -- (1.000,2.000) -- (1.500,3.000) -- cycle;
\filldraw[cefill] (1.000,2.000) -- (2.000,2.000) -- (1.500,3.000) -- cycle;
\node[elemlabel] at (2.167,2.667) {$c_{0}$};
\node[elemlabel] at (2.500,1.500) {$c_{1}$};
\node[elemlabel] at (1.500,0.500) {$c_{2}$};
\node[elemlabel] at (0.500,1.500) {$c_{3}$};
\node[elemlabel] at (0.833,2.667) {$c_{4}$};
\node[elemlabel] at (1.500,2.333) {$c_{5}$};
\draw[hearrow, black!70] (2.134,2.308) -- (2.634,2.808);
\node[helabel, text=black!70] at (2.279,2.610) {$0$};
\draw[hearrow, black!70] (2.593,2.874) -- (1.843,2.874);
\node[helabel, text=black!70] at (2.190,2.760) {$1$};
\draw[hearrow, black!70] (1.746,2.798) -- (1.996,2.298);
\node[helabel, text=black!70] at (1.979,2.592) {$2$};
\draw[hearrow, black!70] (2.750,2.620) -- (2.250,2.120);
\node[helabel, text=black!70] at (2.500,2.253) {$3$};
\draw[hearrow, black!70] (2.130,1.750) -- (2.130,1.250);
\node[helabel, text=black!70] at (2.247,1.500) {$4$};
\draw[hearrow, black!70] (2.250,0.880) -- (2.750,0.380);
\node[helabel, text=black!70] at (2.500,0.747) {$5$};
\draw[hearrow, black!70] (2.870,0.750) -- (2.870,2.250);
\node[helabel, text=black!70] at (2.753,1.500) {$6$};
\draw[hearrow, black!70] (2.620,0.250) -- (2.120,0.750);
\node[helabel, text=black!70] at (2.253,0.500) {$7$};
\draw[hearrow, black!70] (1.750,0.870) -- (1.250,0.870);
\node[helabel, text=black!70] at (1.500,0.753) {$8$};
\draw[hearrow, black!70] (0.880,0.750) -- (0.380,0.250);
\node[helabel, text=black!70] at (0.747,0.500) {$9$};
\draw[hearrow, black!70] (0.750,0.130) -- (2.250,0.130);
\node[helabel, text=black!70] at (1.500,0.247) {$10$};
\draw[hearrow, black!70] (0.250,0.380) -- (0.750,0.880);
\node[helabel, text=black!70] at (0.500,0.747) {$11$};
\draw[hearrow, black!70] (0.870,1.250) -- (0.870,1.750);
\node[helabel, text=black!70] at (0.753,1.500) {$12$};
\draw[hearrow, black!70] (0.750,2.120) -- (0.250,2.620);
\node[helabel, text=black!70] at (0.500,2.253) {$13$};
\draw[hearrow, black!70] (0.130,2.250) -- (0.130,0.750);
\node[helabel, text=black!70] at (0.247,1.500) {$14$};
\draw[hearrow, black!70] (0.366,2.808) -- (0.866,2.308);
\node[helabel, text=black!70] at (0.721,2.610) {$15$};
\draw[hearrow, black!70] (1.004,2.298) -- (1.254,2.798);
\node[helabel, text=black!70] at (1.021,2.592) {$16$};
\draw[hearrow, black!70] (1.157,2.874) -- (0.407,2.874);
\node[helabel, text=black!70] at (0.810,2.760) {$17$};
\draw[hearrow, black!70] (1.300,2.070) -- (1.700,2.070);
\node[helabel, text=black!70] at (1.500,2.161) {$18$};
\draw[hearrow, black!70] (1.792,2.261) -- (1.592,2.661);
\node[helabel, text=black!70] at (1.616,2.411) {$19$};
\draw[hearrow, black!70] (1.408,2.661) -- (1.208,2.261);
\node[helabel, text=black!70] at (1.384,2.411) {$20$};
\node[vreg, minimum size=0.29cm] at (3.000,3.000) {$0$};
\node[vreg, minimum size=0.29cm] at (0.000,3.000) {$1$};
\node[vreg, minimum size=0.29cm] at (0.000,0.000) {$2$};
\node[vreg, minimum size=0.29cm] at (3.000,0.000) {$3$};
\node[vreg, minimum size=0.29cm] at (2.000,2.000) {$4$};
\node[vreg, minimum size=0.29cm] at (1.000,2.000) {$5$};
\node[vreg, minimum size=0.29cm] at (1.000,1.000) {$6$};
\node[vreg, minimum size=0.29cm] at (2.000,1.000) {$7$};
\node[vreg, minimum size=0.29cm] at (1.500,3.000) {$8$};

\end{tikzpicture}

\par\medskip
{\footnotesize\setlength{\tabcolsep}{2.2pt}\begin{tabular}{lccccccccccccccccccccc}
\textit{he} & 0 & 1 & 2 & 3 & 4 & 5 & 6 & 7 & 8 & 9 & 10 & 11 & 12 & 13 & 14 & 15 & 16 & 17 & 18 & 19 & 20 \\ \hline
$V[\,]$ & 4 & 0 & 8 & 0 & 4 & 7 & 3 & 3 & 7 & 6 & 2 & 2 & 6 & 5 & 1 & 1 & 5 & 8 & 5 & 4 & 8 \\
$O[\,]$ & 3 & -- & 19 & 0 & -- & 7 & -- & 5 & -- & 11 & -- & 9 & -- & 15 & -- & 13 & 20 & -- & -- & 2 & 16 \\
\end{tabular}
\par\medskip $Of[\,] = [0,\,3,\,7,\,11,\,15,\,18,\,21]$
}
\par\smallskip
{\small $\V[h]$: vértice de origem de $h$;\quad $\Op[h]$: oposta de $h$ (``--'' = \BORDER);\quad $Of[c]$: primeira
half-edge do elemento $c$.}
\caption{A malha de \code{meshIBHM.txt} com o índice de cada half-edge e os vetores da IBHM tal como construídos pela
classe \code{IBHM}. O último elemento ($c_5$) começa num vértice diferente do usado no artigo, por isso suas três
half-edges têm outra numeração; o restante coincide com a Figura~4 de~\cite{ibhm} subtraindo-se 1 dos índices. O furo
central é delimitado por half-edges de borda.}
\label{fig:ibhm}
\end{figure}

\subsection{A classe \code{IBHM} do repositório}
\label{sec:impl-estruturas}

A classe \code{IBHM} (\code{core/TopologyDataStructures/IBHM.\{h,cpp\}}) segue as ideias acima, com algumas escolhas
próprias que importam para o algoritmo de inserção (\cref{tab:estruturas}). As coordenadas são guardadas em
\code{float}, e as malhas são lidas no formato de \code{malha2Hyb.txt}: número de vértices e de elementos, coordenadas e,
para cada elemento, o número de vértices seguido dos índices.

\begin{table}[ht]
\centering\small
\begin{tabularx}{\linewidth}{>{\raggedright}p{3.4cm}>{\raggedright}X>{\raggedright\arraybackslash}X}
\toprule
Aspecto & Artigo da IBHM~\cite{ibhm} & Classe \code{IBHM}\\
\midrule
Índices & a partir de 1; $\Op[h]=0$ na borda & a partir de 0; $\Op[h]=\BORDER=-1$ (como \code{unsigned}: \code{0xFFFFFFFF})\\
Estado extra de $\Op$ & --- & $\COLLAPSED=-2$ marca uma aresta \emph{colapsada} (de comprimento zero) de um
elemento de interface\\
Offsets & $n_c$ entradas & $n_c+1$ entradas (sentinela final), de modo que $n_v(c)=Of[c+1]-Of[c]$\\
$\mathrm{elem}(h)$ & busca binária & busca binária com \code{std::upper\_bound} nas $n_c+1$ primeiras entradas de
$Of[\,]$ (\code{offsetIterator}); $\nxt$ e $\prv$ também fazem essa busca\\
Construção de $\Op$ & mapa indexado pelo par ordenado de vértices & para cada vértice, lista das half-edges que partem
dele; procura-se a half-edge reversa (\code{buildOppositesLinear})\\
Vetor $VH[\,]$ & sim & \textbf{não}; o teste de borda (\code{isBorder}) gira em torno do vértice\\
Crescimento da malha & não previsto (``limitação'' declarada) & \code{reserveSpaceForElements} aumenta os vetores de
elementos (\code{resize}; as posições novas de $\Op[\,]$ começam com \BORDER) e \code{commitElement} valida cada novo
elemento; \code{numberOfHalfEdges} devolve $Of[n_c]$, o número de half-edges válidas (sem o espaço reservado);
\code{reserveSpaceForNodes} reserva capacidade exata para os pontos (\code{reserve}); \code{addPoint} acrescenta um ponto
com as coordenadas dadas num \code{std::vector} (vazio = origem) e \code{duplicatePoint} acrescenta a cópia de um vértice
existente, com as mesmas coordenadas\\
\bottomrule
\end{tabularx}
\caption{Diferenças entre a IBHM do artigo e a classe \code{IBHM} do repositório.}
\label{tab:estruturas}
\end{table}

Dois pontos merecem destaque. Primeiro, o estado \COLLAPSED{} é uma extensão criada para o algoritmo de inserção:
ele permite que um elemento de interface tenha um lado de comprimento zero sem vizinho, e as consultas que giram em torno
de vértices (\code{isBorder}) sabem pular esse lado. Segundo, a combinação ``contar $\to$ reservar $\to$ validar'' é a
resposta da implementação à limitação apontada no artigo para alterações dinâmicas: o algoritmo de inserção calcula
antes quantos vértices e elementos novos serão necessários (\cref{sec:contagem}), reserva o espaço de uma vez e só
então altera a malha, sem realocações durante as operações.

A interface usada pelo construtor de elementos de interface é pequena: \code{heVertexIndex} ($\V$), \code{heOpposite}
($\Op$), \code{heNext}, \code{hePrevious}, \code{heElement}, \code{isBorder}, \code{setElementVertex} (altera $\V[h]$),
\code{setOpposite} (altera $\Op$ nos dois sentidos quando a oposta é uma half-edge válida, ou em um só sentido quando é
\BORDER{} ou \COLLAPSED), \code{duplicatePoint}, \code{numberOfHalfEdges}, \code{nextAvailableHE} e
\code{commitElement}. Nenhuma coordenada entra nas decisões: o algoritmo é puramente topológico, e as cópias criadas por
\code{duplicatePoint} apenas herdam as coordenadas do vértice de origem.

\subsection{Os vetores da IBHM durante uma inserção}
\label{sec:ibhm-insercao}

A \cref{fig:insercao-vetores} mostra, numa malha híbrida pequena, o que acontece com os vetores quando uma aresta é
aberta. A aresta $4$--$7$ é compartilhada pelos triângulos que contêm as half-edges $a=9$ ($4\to7$) e $b=19$ ($7\to4$);
o vértice 4 é interno e o vértice 7 está na borda.
\begin{enumerate}
\item \textbf{Reserva.} Antes de qualquer alteração, \code{reserveSpaceForElements(4, 1)} acrescenta 4 posições a
$\V[\,]$ e a $\Op[\,]$ (estas com \BORDER) e uma entrada a $Of[\,]$, e \code{reserveSpaceForNodes(1)} reserva capacidade
exata para um vértice, sem mudar o tamanho do vetor de coordenadas (a pré-contagem previu um vértice novo, para a
borda).
\item \textbf{Validação antecipada.} \code{nextAvailableHE()} devolve $Of[n_c]=20$, a primeira half-edge livre, e
\code{commitElement(4)} já declara o elemento 6 válido, gravando $Of[7]=24$. Na IBHM isso precisa acontecer \emph{antes}
dos operadores, porque $\nxt$ e $\prv$ das half-edges 20--23 dependem do offset do elemento.
\item \textbf{Operador em 4} (Canônico, interno): $\V[23]=\V[20]=4$ e $\Op[23]=\COLLAPSED$ (marcado como \textsc{c}).
\item \textbf{Operador em 7} (Canônico, borda): cria o vértice 9, cópia de 7 (\code{duplicatePoint(7)}, mesmas
coordenadas); $\V[21]=9$; a rotação a partir de $b$ troca $\V[19]$ de 7 para 9 e para logo na borda; $\V[22]=7$ e
$\Op[21]=\BORDER$.
\item \textbf{Arestas do meio}: $\Op[20]\leftrightarrow19$ e $\Op[22]\leftrightarrow9$. As antigas opostas $\Op[9]=19$ e
$\Op[19]=9$ são sobrescritas.
\end{enumerate}
Ao todo, uma inserção altera poucas posições dos vetores antigos (em vermelho na figura) e preenche as quatro posições
novas (em laranja). As posições reservadas de $\Op[\,]$ começam com \BORDER: até a v0.3 começavam com 0, um índice
válido de half-edge, e um índice do espaço reservado passado como entrada era aceito como aresta (defeito I3, corrigido;
\cref{sec:defeitos}). No fluxo normal todas são sobrescritas pelos operadores e pela ligação das arestas do meio.

\begin{figure}[ht]
\centering
\begin{subfigure}[b]{0.48\linewidth}\centering\begin{tikzpicture}[scale=2.100, every node/.style={font=\scriptsize}]
\filldraw[cefill] (0.000,0.000) -- (1.000,0.000) -- (1.000,1.000) -- (0.000,1.000) -- cycle;
\filldraw[cefill] (1.000,0.000) -- (2.000,0.000) -- (2.000,1.000) -- (1.000,1.000) -- cycle;
\filldraw[cefill] (0.000,1.000) -- (1.000,1.000) -- (1.000,2.000) -- cycle;
\filldraw[cefill] (0.000,1.000) -- (1.000,2.000) -- (0.000,2.000) -- cycle;
\filldraw[cefill] (1.000,1.000) -- (2.000,1.000) -- (2.000,2.000) -- cycle;
\filldraw[cefill] (1.000,1.000) -- (2.000,2.000) -- (1.000,2.000) -- cycle;
\draw[seledge] (1.000,1.000) -- (1.000,2.000);
\draw[hearrow, black!70] (0.250,0.110) -- (0.750,0.110);
\node[helabel, text=black!70] at (0.500,0.198) {$0$};
\draw[hearrow, black!70] (0.890,0.250) -- (0.890,0.750);
\node[helabel, text=black!70] at (0.802,0.500) {$1$};
\draw[hearrow, black!70] (0.750,0.890) -- (0.250,0.890);
\node[helabel, text=black!70] at (0.500,0.802) {$2$};
\draw[hearrow, black!70] (0.110,0.750) -- (0.110,0.250);
\node[helabel, text=black!70] at (0.198,0.500) {$3$};
\draw[hearrow, black!70] (1.250,0.110) -- (1.750,0.110);
\node[helabel, text=black!70] at (1.500,0.198) {$4$};
\draw[hearrow, black!70] (1.890,0.250) -- (1.890,0.750);
\node[helabel, text=black!70] at (1.802,0.500) {$5$};
\draw[hearrow, black!70] (1.750,0.890) -- (1.250,0.890);
\node[helabel, text=black!70] at (1.500,0.802) {$6$};
\draw[hearrow, black!70] (1.110,0.750) -- (1.110,0.250);
\node[helabel, text=black!70] at (1.198,0.500) {$7$};
\draw[hearrow, black!70] (0.299,1.098) -- (0.799,1.098);
\node[helabel, text=black!70] at (0.589,1.177) {$8$};
\draw[hearrow, red!70!black] (0.902,1.201) -- (0.902,1.701);
\node[helabel, text=red!70!black] at (0.823,1.411) {$9$};
\draw[hearrow, black!70] (0.825,1.675) -- (0.325,1.175);
\node[helabel, text=black!70] at (0.635,1.365) {$10$};
\draw[hearrow, black!70] (0.175,1.325) -- (0.675,1.825);
\node[helabel, text=black!70] at (0.365,1.635) {$11$};
\draw[hearrow, black!70] (0.701,1.902) -- (0.201,1.902);
\node[helabel, text=black!70] at (0.411,1.823) {$12$};
\draw[hearrow, black!70] (0.098,1.799) -- (0.098,1.299);
\node[helabel, text=black!70] at (0.177,1.589) {$13$};
\draw[hearrow, black!70] (1.299,1.098) -- (1.799,1.098);
\node[helabel, text=black!70] at (1.589,1.177) {$14$};
\draw[hearrow, black!70] (1.902,1.201) -- (1.902,1.701);
\node[helabel, text=black!70] at (1.823,1.411) {$15$};
\draw[hearrow, black!70] (1.825,1.675) -- (1.325,1.175);
\node[helabel, text=black!70] at (1.635,1.365) {$16$};
\draw[hearrow, black!70] (1.175,1.325) -- (1.675,1.825);
\node[helabel, text=black!70] at (1.365,1.635) {$17$};
\draw[hearrow, black!70] (1.701,1.902) -- (1.201,1.902);
\node[helabel, text=black!70] at (1.411,1.823) {$18$};
\draw[hearrow, red!70!black] (1.098,1.799) -- (1.098,1.299);
\node[helabel, text=red!70!black] at (1.177,1.589) {$19$};
\node[vreg, minimum size=0.26cm] at (0.000,0.000) {$0$};
\node[vreg, minimum size=0.26cm] at (1.000,0.000) {$1$};
\node[vreg, minimum size=0.26cm] at (2.000,0.000) {$2$};
\node[vreg, minimum size=0.26cm] at (0.000,1.000) {$3$};
\node[vreg, minimum size=0.26cm] at (1.000,1.000) {$4$};
\node[vreg, minimum size=0.26cm] at (2.000,1.000) {$5$};
\node[vreg, minimum size=0.26cm] at (0.000,2.000) {$6$};
\node[vreg, minimum size=0.26cm] at (1.000,2.000) {$7$};
\node[vreg, minimum size=0.26cm] at (2.000,2.000) {$8$};

\end{tikzpicture}
\caption{Antes: a aresta $4$--$7$.}\end{subfigure}\hfill
\begin{subfigure}[b]{0.48\linewidth}\centering\begin{tikzpicture}[scale=2.100, every node/.style={font=\scriptsize}]
\filldraw[cefill] (0.000,0.000) -- (1.000,0.000) -- (1.000,1.000) -- (0.000,1.000) -- cycle;
\filldraw[cefill] (1.000,0.000) -- (2.000,0.000) -- (2.000,1.000) -- (1.000,1.000) -- cycle;
\filldraw[cefill] (0.000,1.000) -- (1.000,1.000) -- (0.859,1.859) -- cycle;
\filldraw[cefill] (0.000,1.000) -- (0.859,1.859) -- (0.000,2.000) -- cycle;
\filldraw[cefill] (1.000,1.000) -- (2.000,1.000) -- (2.000,2.000) -- cycle;
\filldraw[cefill] (1.000,1.000) -- (2.000,2.000) -- (1.141,1.859) -- cycle;
\filldraw[iefill] (1.052,1.085) -- (1.141,1.859) -- (0.859,1.859) -- (0.920,1.060) -- cycle;
\draw[red, line width=1.1pt] (0.920,1.060) -- (1.052,1.085);
\draw[hearrow, black!70] (0.250,0.110) -- (0.750,0.110);
\node[helabel, text=black!70] at (0.500,0.198) {$0$};
\draw[hearrow, black!70] (0.890,0.250) -- (0.890,0.750);
\node[helabel, text=black!70] at (0.802,0.500) {$1$};
\draw[hearrow, black!70] (0.750,0.890) -- (0.250,0.890);
\node[helabel, text=black!70] at (0.500,0.802) {$2$};
\draw[hearrow, black!70] (0.110,0.750) -- (0.110,0.250);
\node[helabel, text=black!70] at (0.198,0.500) {$3$};
\draw[hearrow, black!70] (1.250,0.110) -- (1.750,0.110);
\node[helabel, text=black!70] at (1.500,0.198) {$4$};
\draw[hearrow, black!70] (1.890,0.250) -- (1.890,0.750);
\node[helabel, text=black!70] at (1.802,0.500) {$5$};
\draw[hearrow, black!70] (1.750,0.890) -- (1.250,0.890);
\node[helabel, text=black!70] at (1.500,0.802) {$6$};
\draw[hearrow, black!70] (1.110,0.750) -- (1.110,0.250);
\node[helabel, text=black!70] at (1.198,0.500) {$7$};
\draw[hearrow, black!70] (0.292,1.102) -- (0.792,1.102);
\node[helabel, text=black!70] at (0.576,1.183) {$8$};
\draw[hearrow, black!70] (0.865,1.169) -- (0.794,1.598);
\node[helabel, text=black!70] at (0.750,1.346) {$9$};
\draw[hearrow, black!70] (0.730,1.580) -- (0.300,1.150);
\node[helabel, text=black!70] at (0.583,1.313) {$10$};
\draw[hearrow, black!70] (0.150,1.300) -- (0.580,1.730);
\node[helabel, text=black!70] at (0.313,1.583) {$11$};
\draw[hearrow, black!70] (0.598,1.794) -- (0.169,1.865);
\node[helabel, text=black!70] at (0.346,1.750) {$12$};
\draw[hearrow, black!70] (0.102,1.792) -- (0.102,1.292);
\node[helabel, text=black!70] at (0.183,1.576) {$13$};
\draw[hearrow, black!70] (1.299,1.098) -- (1.799,1.098);
\node[helabel, text=black!70] at (1.589,1.177) {$14$};
\draw[hearrow, black!70] (1.902,1.201) -- (1.902,1.701);
\node[helabel, text=black!70] at (1.823,1.411) {$15$};
\draw[hearrow, black!70] (1.825,1.675) -- (1.325,1.175);
\node[helabel, text=black!70] at (1.635,1.365) {$16$};
\draw[hearrow, black!70] (1.196,1.304) -- (1.696,1.804);
\node[helabel, text=black!70] at (1.403,1.597) {$17$};
\draw[hearrow, black!70] (1.728,1.871) -- (1.299,1.800);
\node[helabel, text=black!70] at (1.467,1.761) {$18$};
\draw[hearrow, black!70] (1.200,1.701) -- (1.129,1.272);
\node[helabel, text=black!70] at (1.239,1.533) {$19$};
\draw[hearrow, orange!50!black] (1.132,1.255) -- (1.182,1.688);
\node[helabel, text=orange!50!black] at (1.337,1.222) {$20$};
\draw[hearrow, orange!50!black] (0.812,1.683) -- (0.847,1.236);
\node[helabel, text=orange!50!black] at (0.629,1.709) {$22$};
\node[vreg, minimum size=0.26cm] at (0.000,0.000) {$0$};
\node[vreg, minimum size=0.26cm] at (1.000,0.000) {$1$};
\node[vreg, minimum size=0.26cm] at (2.000,0.000) {$2$};
\node[vreg, minimum size=0.26cm] at (0.000,1.000) {$3$};
\node[vcol, minimum size=0.26cm] at (1.000,1.000) {$4$};
\node[vreg, minimum size=0.26cm] at (2.000,1.000) {$5$};
\node[vreg, minimum size=0.26cm] at (0.000,2.000) {$6$};
\node[vie, minimum size=0.26cm] at (0.859,1.859) {$7$};
\node[vreg, minimum size=0.26cm] at (2.000,2.000) {$8$};
\node[vie, minimum size=0.26cm] at (1.141,1.859) {$9$};

\end{tikzpicture}
\caption{Depois: elemento $6=[4,9,7,4]$.}\end{subfigure}
\par\medskip
{\scriptsize\setlength{\tabcolsep}{2.1pt}\begin{tabular}{lcccccccccccccccccccccccc}
\textit{he} & 0 & 1 & 2 & 3 & 4 & 5 & 6 & 7 & 8 & 9 & 10 & 11 & 12 & 13 & 14 & 15 & 16 & 17 & 18 & 19 & 20 & 21 & 22 & 23 \\ \hline
$\V$ antes & 0 & 1 & 4 & 3 & 1 & 2 & 5 & 4 & 3 & 4 & 7 & 3 & 7 & 6 & 4 & 5 & 8 & 4 & 8 & 7 &  &  &  &  \\
$\V$ depois & 0 & 1 & 4 & 3 & 1 & 2 & 5 & 4 & 3 & 4 & 7 & 3 & 7 & 6 & 4 & 5 & 8 & 4 & 8 & \textcolor{red!75!black}{\textbf{9}} & \textcolor{orange!50!black}{\textbf{4}} & \textcolor{orange!50!black}{\textbf{9}} & \textcolor{orange!50!black}{\textbf{7}} & \textcolor{orange!50!black}{\textbf{4}} \\ \hline
$\Op$ antes & -- & 7 & 8 & -- & -- & -- & 14 & 1 & 2 & 19 & 11 & 10 & -- & -- & 6 & -- & 17 & 16 & -- & 9 &  &  &  &  \\
$\Op$ depois & -- & 7 & 8 & -- & -- & -- & 14 & 1 & 2 & \textcolor{red!75!black}{\textbf{22}} & 11 & 10 & -- & -- & 6 & -- & 17 & 16 & -- & \textcolor{red!75!black}{\textbf{20}} & \textcolor{orange!50!black}{\textbf{19}} & \textcolor{orange!50!black}{\textbf{--}} & \textcolor{orange!50!black}{\textbf{9}} & \textcolor{orange!50!black}{\textbf{\textsc{c}}} \\
\end{tabular}
\par\smallskip $Of$ antes $= [0,\,4,\,8,\,11,\,14,\,17,\,20]$\qquad $Of$ depois $= [0,\,4,\,8,\,11,\,14,\,17,\,20,\,24]$
}
\caption{Uma inserção vista pelos vetores da IBHM. Na tabela, vermelho indica entradas antigas que mudaram e laranja as
posições do novo elemento (half-edges 20--23); ``--'' é \BORDER{} e \textsc{c} é \COLLAPSED. O novo elemento é
$[F,G,H,E]=[4,9,7,4]$: colapsado em 4 (interno) e aberto em 7 (borda), onde surgiu o vértice 9.}
\label{fig:insercao-vetores}
\end{figure}
